% Macros for content
\newcommand{\systemname}{Ihr System}
\newcommand{\teamname}{Ihr Team}
\newcommand{\workshopdate}{\today}
\newcommand{\iteration}{v1.0}

\newcommand{\businesscaseheadline}{Business Case}
\newcommand{\businesscasecontent}{
Kurze Beschreibung des Business Case oder ökonomischen Treibers hinter dem Softwaresystem.
}
\newcommand{\businesscasehint}{
Wenn das Unternehmen Geld verdienen muss, gibt es einen wirtschaftlichen Treiber und einen Business Case hinter jeder Initiative. Wenn nicht, ist es ein Hobby oder ein NGO-Projekt.

Am Anfang steht immer das WARUM. -- Beginnen Sie mit dem WARUM.

Eine häufige Falle ist, dass das Entwicklungsteam oder schlimmer noch das gesamte Produktteam sich der spezifischen wirtschaftlichen Treiber und des Business Case hinter diesen Software-Initiativen nicht bewusst ist.
}

\newcommand{\functionaloverviewheadline}{Functional Overview}
\newcommand{\functionaloverviewcontent}{
Die wichtigsten funktionalen Anforderungen auf hoher Ebene.
}
\newcommand{\functionaloverviewhint}{
Wenn Sie als Produktteam die funktionale Übersicht der Initiative erarbeiten, stellen Sie sich eine Produktbox Ihrer Lösung vor - was steht auf Ihrer Box?

Was steht auf Ihrer Software-Produktbox?

Denken Sie die Produktbox, um Antworten auf folgende Fragen zu erhalten:

\begin{itemize}
\item Welchen geschäftlichen Wert liefert die Software?
\item Was sind die übergeordneten Geschäftsfunktionen?
\end{itemize}

Mit den Antworten auf diese Fragen können Sie die wichtigsten funktionalen Anforderungen auf hoher Ebene beschreiben.
}

\newcommand{\qualitygoalsheadline}{Quality Goals}
\newcommand{\qualitygoalscontent}{
Die drei wichtigsten Qualitätsziele für die Architektur mit höchster Priorität für die wichtigsten Stakeholder.
}
\newcommand{\qualitygoalshint}{
Hier arbeiten Sie die drei wichtigsten Qualitätsziele der Stakeholder für das Softwaresystem heraus.

Es ist wichtig, dass diese Qualitätsziele jedem Teammitglied bekannt sind, da sie Ihre Architektur prägen.

Der Standard \href{https://iso25000.com/index.php/en/iso-25000-standards/iso-25010}{ISO/IEC 25010} bietet einen Überblick über mögliche interessante Qualitätsattribute, das moderne arc42-Qualitätsmodell (\href{https://quality.arc42.org/}{Q42}) sammelt mehr als 130 Qualitätsattribute mit ausführlichen Erklärungen und Beispielen.
}

\newcommand{\businesscontextheadline}{Business Context}
\newcommand{\businesscontextcontent}{
Trennen Sie Ihr System als Black Box von allen Kommunikationspartnern. Kommunikationspartner sind benachbarte externe Systeme und Benutzer.
}
\newcommand{\businesscontexthint}{
Der Business Context konzentriert sich auf Ihr System als Black Box. Die Aufmerksamkeit richtet sich auf die Systemumgebung und wie das System mit allen externen Systemen darin interagiert.

Mit dem Business Context können Sie die Kommunikationspfade und damit potenzielle technische Risiken identifizieren.

Es ist ebenso eine gute Möglichkeit, mit verschiedenen Stakeholdern den Umfang des Systems und die Kommunikationspfade zu verschiedenen externen Systemen zu diskutieren.

Wenn Sie versuchen, den Business Context herauszuarbeiten, versuchen Sie die Benutzer Ihres Systems und alle benachbarten Systeme zu identifizieren.
}

\newcommand{\orgconstraintsheadline}{Organisational Constraints}
\newcommand{\orgconstraintscontent}{
Jede organisatorische Anforderung, die die Entscheidungsfreiheit der Software-Architekten einschränkt.
}
\newcommand{\orgconstraintshint}{
Randbedingungen bestimmen den Grad Ihrer Freiheit, als Engineering-Team Entscheidungen zu treffen.

In diesem Teil arbeiten Sie alle Anforderungen heraus, die Ihre Freiheit als Produktteam bei der Entwicklung des (Software-)Produkts einschränken.

Organisatorische Randbedingungen beinhalten oft Zeit und Geld. Wir haben auch gesehen, dass es Einschränkungen bei den verwendeten Methoden (wie SAFe, Scrum, …) oder Techniken gibt.
}

\newcommand{\techconstraintsheadline}{Technical Constraints}
\newcommand{\techconstraintscontent}{
Jede technische Anforderung, die die Entscheidungsfreiheit der Software-Architekten einschränkt.
}
\newcommand{\techconstraintshint}{
Randbedingungen bestimmen den Grad Ihrer Freiheit, als Engineering-Team Entscheidungen zu treffen.

In diesem Teil arbeiten Sie alle Anforderungen heraus, die Ihre Freiheit als Produktteam bei der Entwicklung des (Software-)Produkts einschränken.

In größeren Unternehmen gibt es oft viele technische Randbedingungen. Es gibt oft ein übergeordnetes Architekturkomitee, das die Standards für Technologieentscheidungen setzt.

Wenn Sie sich beispielsweise entscheiden, eine Single-Page-Anwendung zu bauen, wird das übergeordnete Architekturkomitee Angular als SPA-Framework festlegen, sodass es Expertisesynergien zwischen Teams gibt.
}

\newcommand{\archhypothesesheadline}{Architectural Hypotheses}
\newcommand{\archhypothesescontent}{
Resultierende Architekturhypothesen und wichtige, teure, großskalige oder riskante Entscheidungen inkl. Begründungen.
}
\newcommand{\archhypotheseshint}{
Herzlichen Glückwunsch - Sie haben Ihren architektonischen Spielplatz erstellt.

Jetzt können Sie Ihre ersten Architekturhypothesen basierend auf Ihrem aktuellen Teamwissen erstellen.

Beachten Sie: Es geht hier um übergeordnete architektonische Annahmen, nicht um Mikro-Entscheidungen wie Designfragen.

Idealerweise können Sie Ihre architektonischen Annahmen in Entwürfe für Architecture Decision Records umwandeln.
}

\newcommand{\challengesrisksheadline}{Technical Challenges \& Risks}
\newcommand{\challengesriskscontent}{
Identifizierte aktuell bekannte Herausforderungen und technische Risiken.
}
\newcommand{\challengesriskshint}{
Basierend auf dem Spielplatz, den Sie erstellt haben, wie Business Context und Scope, Qualitätsziele und Randbedingungen, können Sie die bevorstehenden Herausforderungen und die daraus resultierenden Risiken erfassen.

Es ist sehr wichtig, dies herauszuarbeiten, kontinuierlich zu bewerten und den wichtigsten Stakeholdern im System zu kommunizieren.

Dieser Abschnitt beantwortet die Frage, wie Sie die aktuelle Situation im Verhältnis zum erarbeiteten architektonischen Spielplatz bewerten.
}
